% SPDX-License-Identifier: GPL-3.0-or-later
\chapter{Laboratorul Orbitar, La Manșa Navei și Misiunile Cadetului}
\label{chap:laborator}

Dincolo de enciclopedia corpurilor cerești, \textbf{Orbitar} oferă unelte interactive de fizică aplicată și explorare: mini-jocul de pilotaj spațial \textbf{„La Manșa Navei (Mod Pilot)”}, \textbf{Laboratorul Orbitar} (calculatoarele de greutate, săritură și vârstă pe alte lumi), \textbf{Rigla cu Viteza Luminii}, jocul de orientare cosmică \textbf{„Misiunile Orbitar \& Jurnalul de Bord”} și \textbf{Normalizatorul Fonetic Român}.

\section{„La Manșa Navei (Mod Pilot)”: Cele 3 Astronave Didactice Românești}

În modul First-Person Cockpit (activat prin tasta \texttt{P} în \textbf{Orbitar}), copiii preiau manșa a \textbf{trei nave spațiale didactice} care omagiază pionierii aviației românești și literatura română, ilustrând scara reală a timpului și a distanțelor cosmice:

\begin{tcolorbox}[orbitardata, title={\textbf{Flota Spațială Didactică din Orbitar (\texttt{SpacecraftPilot.ts})}}]
\begin{enumerate}[leftmargin=1.5em, itemsep=5pt]
  \item \textbf{Racheta „Vlaicu-1”, Propulsie Chimică ($50.000$ de kilometri pe oră / $0{,}0046\%\,c$):}\\
  Numită în onoarea inginerului și pionierului aviației românești \textbf{Aurel Vlaicu} (1882--1913). Are viteza caracteristică a sondelor interplanetare actuale (\emph{New Horizons}, \emph{Apollo}, \emph{Voyager}) și urmează \textbf{arcul de transfer orbital Hohmann} și manevrele de asistență gravitațională: călătoria de la Pământ la Lună durează aproximativ $3\text{ zile}$, până la Marte aproximativ $7\text{ luni}$, până la Jupiter aproximativ $5\text{ ani și } 6\text{ luni}$, iar până la Neptun aproximativ $12\text{ ani}$!
  \item \textbf{Nava „Coandă-X”, Propulsie cu Fuziune \& Plasmă ($1.500.000$ de kilometri pe oră / $0{,}139\%\,c$):}\\
  Numită în onoarea savantului și inventatorului avionului cu reacție \textbf{Henri Coandă} (1886--1972). Reprezintă tehnologia spațială a viitorului apropiat, de 30 de ori mai rapidă decât rachetele chimice: ajunge la Lună în aproximativ $3\text{ ore}$, pe Marte în aproximativ $2\text{ zile}$, la Jupiter în aproximativ $18\text{ zile}$ și la Neptun în aproximativ $4\text{ luni}$!
  \item \textbf{Astronava „Luceafărul”, Motor Fotonic ($1c = 1.079.252.849$ kilometri pe oră) \& Salt Hiperspațial:}\\
  Inspirată de poemul cosmogonic \emph{„Luceafărul”} al lui \textbf{Mihai Eminescu} (\emph{„Porni luceafărul. Creșteau / În cer a lui aripe, / Și căi de mii de ani treceau / În tot atâtea clipe”}). În Sistemul Solar călătorește exact cu \textbf{viteza luminii ($1c$)}: $1{,}3\text{ secunde}$ până la Lună, $8\text{ minute și } 20\text{ de secunde}$ de la Soare la Pământ și aproximativ $4\text{ ore}$ până la Neptun. Pentru țintele din Orizonturile 2 și 3 (stele, exoplanete și galaxii), activează \textbf{Motorul Hiperspațial (Warp)} și explică la sosire în ce epocă istorică a plecat lumina de la acea stea sau galaxie!
\end{enumerate}
\end{tcolorbox}

\begin{table}[H]
\centering
\small
\caption{Compararea duratei călătoriilor de la Pământ cu cele 3 astronave din Orbitar}
\resizebox{\textwidth}{!}{%
\begin{tabular}{@{}llrrr@{}}
\toprule
\textbf{Destinație} & \textbf{Distanță tipică} & \textbf{Racheta „Vlaicu-1”} & \textbf{Nava „Coandă-X”} & \textbf{Astronava „Luceafărul”} \\
 & & ($50.000\text{ km/h}$) & ($1.500.000\text{ km/h}$) & (Viteza luminii $1c$) \\
\midrule
Luna & $384.400\text{ km}$ & $3\text{ zile}$ & $3\text{ ore}$ & $1{,}3\text{ secunde}$ \\
Soarele & $1{,}00\text{ UA}$ ($150\text{ mil. km}$) & $\sim 6\text{ luni}$ & $4\text{ zile}$ & $8\text{ min } 20\text{ sec}$ \\
Marte & $\sim 0{,}52\text{ UA}$ (min.) & $7\text{ luni}$ (Hohmann) & $2\text{ zile}$ & $\sim 4{,}3\text{ minute}$ \\
Jupiter & $\sim 4{,}20\text{ UA}$ & $5\text{ ani și } 6\text{ luni}$ & $18\text{ zile}$ & $\sim 35\text{ minute}$ \\
Saturn & $\sim 8{,}58\text{ UA}$ & $7\text{ ani}$ & $36\text{ de zile}$ & $1\text{ oră și } 11\text{ min}$ \\
Uranus & $\sim 18{,}20\text{ UA}$ & $9\text{ ani și } 6\text{ luni}$ & $76\text{ de zile}$ & $2\text{ ore și } 31\text{ min}$ \\
Neptun & $\sim 29{,}07\text{ UA}$ & $12\text{ ani}$ & $4\text{ luni}$ & $4\text{ ore și } 2\text{ min}$ \\
Pluto & $\sim 38{,}5\text{ UA}$ & $14\text{ ani și } 6\text{ luni}$ & $5\text{ luni}$ & $5\text{ ore și } 20\text{ min}$ \\
Proxima Centauri b & $4{,}24\text{ ani-lumină}$ & $\sim 81.000\text{ de ani}$ & $\sim 3.050\text{ de ani}$ & $4\text{ ani și } 3\text{ luni}$ (sau Warp) \\
Galaxia Andromeda & $2{,}54\text{ mil. a.l.}$ & $\sim 50\text{ mld. ani}$ & $\sim 1{,}8\text{ mld. ani}$ & $2{,}54\text{ mil. ani}$ (sau Warp) \\
\bottomrule
\end{tabular}%
}
\end{table}

\section{Laboratorul Orbitar: Cât cântăresc și cât de sus sar pe alte lumi?}

Mulți oameni confundă \textbf{masa} cu \textbf{greutatea}, dar în astronomie diferența dintre ele este esențială:
\begin{itemize}[leftmargin=1.5em, itemsep=2pt]
  \item \textbf{Masa ($m$, măsurată în kilograme)} reprezintă cantitatea de materie din corpul tău și rămâne \textbf{neschimbată} oriunde în Univers.
  \item \textbf{Greutatea ($G = m \cdot g_{\text{local}}$)} reprezintă forța cu care gravitația corpului ceresc te atrage spre sol. În \textbf{Laboratorul Orbitar} (\texttt{OrbitarLab.ts}), greutatea aparentă și înălțimea săriturii pe o lume cu raportul gravitațional $g_{\text{raport}} = g_{\text{local}} / g_\oplus$ se calculează prin relațiile:
  \[
    M_{\text{echivalent}} = M_{\text{Pământ}} \times g_{\text{raport}}, \qquad h_{\text{săritură}} = \frac{h_{\text{Pământ}}}{g_{\text{raport}}}
  \]
\end{itemize}

Tabelul de mai jos prezintă valorile exacte calculate de \textbf{Laboratorul Orbitar} pentru un copil cu masa de \textbf{35 kg} care pe Pământ sare \textbf{40 cm} în înălțime:

\begingroup
\small
\begin{longtable}{@{} >{\raggedright\arraybackslash}p{2.5cm} >{\raggedright\arraybackslash}p{3.1cm} r r r @{}}
\caption{Cât cântărește și cât de sus sare un copil de 35 kg (săritură 40 cm pe Pământ) pe toate cele 50 de corpuri naturale din Sistemul Solar} \\
\toprule
\textbf{Corp ceresc} & \textbf{Categorie} & \textbf{Gravitație ($g_\oplus$)} & \textbf{Cântarul arată} & \textbf{Înălțimea săriturii} \\
\midrule
\endfirsthead
\multicolumn{5}{c}{{\footnotesize\itshape (continuare de pe pagina anterioară)}} \\
\toprule
\textbf{Corp ceresc} & \textbf{Categorie} & \textbf{Gravitație ($g_\oplus$)} & \textbf{Cântarul arată} & \textbf{Înălțimea săriturii} \\
\midrule
\endhead
\bottomrule
\endfoot
Soarele & Stea & 27,940 $\times g_\oplus$ & \textbf{977,9 kg} & \textbf{1,4 cm} \\
Mercur & Planetă Telurică (Stâncoasă) & 0,380 $\times g_\oplus$ & \textbf{13,3 kg} & \textbf{1,05 metri} \\
Venus & Planetă Telurică (Stâncoasă) & 0,900 $\times g_\oplus$ & \textbf{31,5 kg} & \textbf{44,4 cm} \\
Pământul & Planetă Telurică (Stâncoasă) & 1,000 $\times g_\oplus$ & \textbf{35,0 kg} & \textbf{40,0 cm} \\
Marte & Planetă Telurică (Stâncoasă) & 0,380 $\times g_\oplus$ & \textbf{13,3 kg} & \textbf{1,05 metri} \\
Jupiter & Gigant Gazos & 2,530 $\times g_\oplus$ & \textbf{88,5 kg} & \textbf{15,8 cm} \\
Saturn & Gigant Gazos & 1,060 $\times g_\oplus$ & \textbf{37,1 kg} & \textbf{37,7 cm} \\
Uranus & Gigant de Gheață & 0,890 $\times g_\oplus$ & \textbf{31,2 kg} & \textbf{44,9 cm} \\
Neptun & Gigant de Gheață & 1,140 $\times g_\oplus$ & \textbf{39,9 kg} & \textbf{35,1 cm} \\
Luna & Satelit Natural (Lună) & 0,165 $\times g_\oplus$ & \textbf{5,8 kg} & \textbf{2,42 metri} \\
Io & Satelit Natural (Lună) & 0,183 $\times g_\oplus$ & \textbf{6,4 kg} & \textbf{2,19 metri} \\
Europa & Satelit Natural (Lună) & 0,134 $\times g_\oplus$ & \textbf{4,7 kg} & \textbf{2,99 metri} \\
Ganymede & Satelit Natural (Lună) & 0,146 $\times g_\oplus$ & \textbf{5,1 kg} & \textbf{2,74 metri} \\
Callisto & Satelit Natural (Lună) & 0,126 $\times g_\oplus$ & \textbf{4,4 kg} & \textbf{3,17 metri} \\
Phobos & Satelit Natural (Lună) & 0,001 $\times g_\oplus$ & \textbf{0,0 kg} & \textbf{689,66 metri} \\
Deimos & Satelit Natural (Lună) & 0,000 $\times g_\oplus$ & \textbf{0 kg (Imponderabilitate)} & \textbf{Plutești în spațiu ($\infty$)} \\
Mimas & Satelit Natural (Lună) & 0,006 $\times g_\oplus$ & \textbf{0,2 kg} & \textbf{61,54 metri} \\
Enceladus & Satelit Natural (Lună) & 0,011 $\times g_\oplus$ & \textbf{0,4 kg} & \textbf{36,36 metri} \\
Tethys & Satelit Natural (Lună) & 0,015 $\times g_\oplus$ & \textbf{0,5 kg} & \textbf{26,67 metri} \\
Dione & Satelit Natural (Lună) & 0,024 $\times g_\oplus$ & \textbf{0,8 kg} & \textbf{16,67 metri} \\
Rhea & Satelit Natural (Lună) & 0,027 $\times g_\oplus$ & \textbf{0,9 kg} & \textbf{14,81 metri} \\
Titan & Satelit Natural (Lună) & 0,138 $\times g_\oplus$ & \textbf{4,8 kg} & \textbf{2,90 metri} \\
Hyperion & Satelit Natural (Lună) & 0,002 $\times g_\oplus$ & \textbf{0,1 kg} & \textbf{200,00 metri} \\
Iapetus & Satelit Natural (Lună) & 0,023 $\times g_\oplus$ & \textbf{0,8 kg} & \textbf{17,39 metri} \\
Phoebe & Satelit Natural (Lună) & 0,004 $\times g_\oplus$ & \textbf{0,1 kg} & \textbf{100,00 metri} \\
Miranda & Satelit Natural (Lună) & 0,008 $\times g_\oplus$ & \textbf{0,3 kg} & \textbf{50,00 metri} \\
Ariel & Satelit Natural (Lună) & 0,025 $\times g_\oplus$ & \textbf{0,9 kg} & \textbf{16,00 metri} \\
Umbriel & Satelit Natural (Lună) & 0,023 $\times g_\oplus$ & \textbf{0,8 kg} & \textbf{17,39 metri} \\
Titania & Satelit Natural (Lună) & 0,039 $\times g_\oplus$ & \textbf{1,4 kg} & \textbf{10,26 metri} \\
Oberon & Satelit Natural (Lună) & 0,035 $\times g_\oplus$ & \textbf{1,2 kg} & \textbf{11,43 metri} \\
Triton & Satelit Natural (Lună) & 0,079 $\times g_\oplus$ & \textbf{2,8 kg} & \textbf{5,06 metri} \\
Charon & Satelit Natural (Lună) & 0,029 $\times g_\oplus$ & \textbf{1,0 kg} & \textbf{13,79 metri} \\
Pluto & Planetă Pitică & 0,063 $\times g_\oplus$ & \textbf{2,2 kg} & \textbf{6,35 metri} \\
Ceres & Planetă Pitică & 0,029 $\times g_\oplus$ & \textbf{1,0 kg} & \textbf{13,79 metri} \\
Eris & Planetă Pitică & 0,084 $\times g_\oplus$ & \textbf{2,9 kg} & \textbf{4,76 metri} \\
Haumea & Planetă Pitică & 0,041 $\times g_\oplus$ & \textbf{1,4 kg} & \textbf{9,76 metri} \\
Makemake & Planetă Pitică & 0,050 $\times g_\oplus$ & \textbf{1,8 kg} & \textbf{8,00 metri} \\
4 Vesta & Asteroid & 0,025 $\times g_\oplus$ & \textbf{0,9 kg} & \textbf{16,00 metri} \\
2 Pallas & Asteroid & 0,022 $\times g_\oplus$ & \textbf{0,8 kg} & \textbf{18,18 metri} \\
101955 Bennu & Asteroid & 0,000 $\times g_\oplus$ & \textbf{0 kg (Imponderabilitate)} & \textbf{Plutești în spațiu ($\infty$)} \\
99942 Apophis & Asteroid & 0,000 $\times g_\oplus$ & \textbf{0 kg (Imponderabilitate)} & \textbf{Plutești în spațiu ($\infty$)} \\
16 Psyche & Asteroid & 0,015 $\times g_\oplus$ & \textbf{0,5 kg} & \textbf{26,67 metri} \\
\end{longtable}
\endgroup

\section{Laboratorul Orbitar: Câți ani și câte zile am pe alte planete?}

Pe Pământ, un \textbf{an} durează $365{,}256\text{ zile}$ (timpul unei rotații complete în jurul Soarelui), iar o \textbf{zi} durează $24\text{ de ore}$ (o rotație în jurul propriei axe). Pentru un copil în vârstă de \textbf{10 ani pământești} ($3.652{,}56\text{ zile pământești} = 87.661{,}4\text{ ore}$), iată vârsta exprimată în ani și zile locale pe fiecare planetă și planetă pitică:

\begin{table}[H]
\centering
\small
\caption{Vârsta unui copil de 10 ani pământești pe planetele și planetele pitice din Orbitar}
\begin{tabular}{@{}lrrrr@{}}
\toprule
\textbf{Lumea} & \textbf{Durata unui an} & \textbf{Durata unei zile} & \textbf{Vârsta în ani locali} & \textbf{Vârsta în zile locale} \\
\midrule
Mercur & 88,0 zile & 1407,6 ore & \textbf{41,5 ani} & \textbf{62 zile locale} \\
Venus & 224,7 zile & 5832,5 ore & \textbf{16,3 ani} & \textbf{15 zile locale} \\
Pământul & 365,3 zile & 23,9 ore & \textbf{10,0 ani} & \textbf{3.663 zile locale} \\
Marte & 687,0 zile & 24,6 ore & \textbf{5,317 ani} & \textbf{3.560 zile locale} \\
Jupiter & 4332,6 zile & 9,9 ore & \textbf{0,843 ani} & \textbf{8.832 zile locale} \\
Saturn & 10759,2 zile & 10,7 ore & \textbf{0,339 ani} & \textbf{8.193 zile locale} \\
Uranus & 30688,5 zile & 17,2 ore & \textbf{0,119 ani} & \textbf{5.085 zile locale} \\
Neptun & 60182,0 zile & 16,1 ore & \textbf{0,061 ani} & \textbf{5.441 zile locale} \\
Pluto & 90560,0 zile & 153,3 ore & \textbf{0,040 ani} & \textbf{572 zile locale} \\
Ceres & 1682,0 zile & 9,1 ore & \textbf{2,172 ani} & \textbf{9.665 zile locale} \\
Eris & 203830,0 zile & 25,9 ore & \textbf{0,018 ani} & \textbf{3.385 zile locale} \\
Haumea & 103774,0 zile & 3,9 ore & \textbf{0,035 ani} & \textbf{22.391 zile locale} \\
Makemake & 111845,0 zile & 22,8 ore & \textbf{0,033 ani} & \textbf{3.840 zile locale} \\
\bottomrule
\end{tabular}
\end{table}

\section{Misiunile Orbitar și Jurnalul de Bord (Vânătoarea de Comori Cosmică)}

Pentru a transforma studiul Sistemului Solar într-o aventură interactivă, modulul \textbf{Misiunile Orbitar} (\texttt{OrbitarMissions.ts}) propune 6 provocări de explorare. Fiecare misiune îi cere copilului să găsească ținta în scena 3D și apoi să răspundă corect la întrebarea grilă pentru a debloca insigna în \textbf{Jurnalul de Bord}:

\begin{tcolorbox}[orbitarquiz, title={\textbf{Cele 6 Misiuni Oficiale și Insigne din Jurnalul de Bord Orbitar}}]
\begin{enumerate}[leftmargin=1.5em, itemsep=4pt]
  \item \textbf{Misiunea 1: Insigna „Ochiul de Aur (JWST)” (țintă: \texttt{jwst}):}\\
  \emph{Provocare:} Găsește în Orbitar telescopul spațial aflat în punctul L2 și dă click pe el!\\
  \emph{Indiciu:} Are o oglindă aurită în formă de fagure și orbitează chiar lângă Pământ.
  \item \textbf{Misiunea 2: Insigna „Explorator Marțian” (țintă: \texttt{mars}):}\\
  \emph{Provocare:} Găsește Planeta Roșie unde au amartizat roverele Perseverance și Curiosity!\\
  \emph{Indiciu:} Este a patra planetă de la Soare, aflată imediat după orbita Pământului.
  \item \textbf{Misiunea 3: Insigna „Stăpânul Inelelor” (țintă: \texttt{saturn}):}\\
  \emph{Provocare:} Descoperă bijuteria Sistemului Solar, gigantul gazos înconjurat de inele strălucitoare de gheață!\\
  \emph{Indiciu:} Este a șasea planetă de la Soare, după uriașul Jupiter.
  \item \textbf{Misiunea 4: Insigna „Comandant Orbital (ISS)” (țintă: \texttt{iss}):}\\
  \emph{Provocare:} Găsește laboratorul spațial locuit de astronauți care zboară pe orbita joasă a Pământului (la aproximativ $420$ de kilometri altitudine)!\\
  \emph{Indiciu:} Apropie-te de Pământ și selectează Stația Spațială Internațională.
  \item \textbf{Misiunea 5: Insigna „Inima din Kuiper” (țintă: \texttt{pluto}):}\\
  \emph{Provocare:} Caută la marginea Sistemului Solar planeta pitică ce poartă pe suprafața ei un ghețar în formă de inimă!\\
  \emph{Indiciu:} Se află în Centura Kuiper, dincolo de orbita albastră a lui Neptun.
  \item \textbf{Misiunea 6: Insigna „Mesager Interstelar” (țintă: \texttt{voyager1}):}\\
  \emph{Provocare:} Găsește sonda lansată în 1977 care poartă în spațiul interstelar Discul de Aur cu salutul în limba română!\\
  \emph{Indiciu:} Este cel mai îndepărtat obiect creat de om, așa că poți căuta sonda \emph{Voyager 1}!
\end{enumerate}
\end{tcolorbox}

\section{Ghid de Pronunție Astronomică în Limba Română}

Multe nume de stele, sateliți, cratere, astronomi și sonde spațiale provin din greaca veche, latină, arabă, hawaiiană sau engleză. Iată ghidul fonetic utilizat de modulul \texttt{RomanianSpeechNormalizer.ts} din \textbf{Orbitar} pentru pronunțarea lor corectă:

\begin{table}[H]
\centering
\small
\caption{Dicționarul fonetic al termenilor astronomici și al misiunilor din Orbitar}
\resizebox{\textwidth}{!}{%
\begin{tabular}{@{}lll@{}}
\toprule
\textbf{Grafie oficială} & \textbf{Pronunție recomandată} & \textbf{Explicație / Origine} \\
\midrule
Betelgeuse & \emph{Be-tel-geuz} & Supergigantă roșie în constelația Orion \\
Enceladus & \emph{En-țe-la-dus} & Satelitul cu gheizere de gheață al lui Saturn \\
Ganymede & \emph{Ga-ni-me-de} & Cel mai mare satelit al lui Jupiter \\
Callisto & \emph{Ca-lis-to} & Satelit galileean exterior al lui Jupiter \\
Phobos / Stickney & \emph{Fo-bos / Stic-ni} & Satelitul interior al lui Marte și marele său crater \\
Charon & \emph{Ha-ron} & Satelitul binar al planetei pitice Pluto \\
Haumea / Makemake & \emph{Ha-u-me-a / Ma-ke-ma-ke} & Planete pitice din Centura Kuiper \\
16 Psyche / 101955 Bennu & \emph{Psi-he / Be-nu} & Asteroidul metalic și asteroidul carbonifer \\
99942 Apophis / 2 Pallas & \emph{A-po-fis / Pa-las} & Asteroizi din Sistemul Solar \\
Centura Kuiper / Norul Oort & \emph{Cai-păr / Ort} & După astronomii olandezi Gerard Kuiper și Jan Oort \\
Clyde Tombaugh & \emph{Claid Tom-bo} & Descoperitorul lui Pluto (1930); \emph{Tombaugh Regio} \\
Christiaan Huygens & \emph{Cris-ti-an Hoi-ghens} & Descoperitorul lui Titan și sonda europeană \emph{Huygens} \\
William Herschel & \emph{Ui-li-am Her-șăl} & Descoperitorul planetei Uranus (1781) \\
67P/Churyumov-Gerasimenko & \emph{67P Ciu-riu-mov Ghe-ra-si-men-co} & Cometa explorată de misiunea \emph{Rosetta \& Philae} \\
James Webb (JWST) & \emph{Geims Ueb} & Telescopul spațial în infraroșu din punctul L2 \\
Hubble / Spitzer / Kepler & \emph{Ha-băl / Șpi-țăr / Che-plăr} & Marile telescoape spațiale NASA \\
Voyager 1 \& 2 & \emph{Vo-ia-ger unu și doi} & Sondele interstelare lansate în 1977 \\
New Horizons & \emph{Niu Ho-rai-zăns} & Sonda care a zburat pe lângă Pluto în 2015 \\
Parker Solar Probe & \emph{Sonda Solară Par-căr} & Sonda care studiază coroana Soarelui \\
Rosetta \& Philae & \emph{Ro-ze-ta și Fi-le} & Sonda orbitală și modulul de asolizare pe cometă \\
JUICE & \emph{Juis} & \emph{Jupiter Icy Moons Explorer} (ESA) \\
Perseverance \& Ingenuity & \emph{Per-se-ve-răns și In-ge-niu-i-ti} & Roverul și elicopterul NASA din Craterul Jezero \\
Curiosity & \emph{Chiu-ri-o-si-ti} & Roverul nuclear NASA din Craterul Gale \\
Sagittarius A* (Sgr A*) & \emph{Sa-gi-ta-ri-us A stea} & Gaura neagră supermasivă din centrul Căii Lactee \\
TRAPPIST-1 & \emph{Tra-pist unu} & Sistemul cu 7 exoplanete stâncoase \\
\bottomrule
\end{tabular}%
}
\end{table}

\clearpage
